\documentclass[11pt,a4paper]{article}
\usepackage{review-style}
% Invented teaching material demonstrating structure, not actual course content.
% Replace these phrases with the current slides' original wording.
% Keep lecture order, source synthesis, and each instance with its parent.
% For proof-oriented courses, also read references/mathematical-review.md.
% Use the course's proof badges; a full displayed proof does not imply that
% independent reconstruction is required. Define symbols and randomness first.
% For substantive proofs, state the Core idea, then named actions and displays.
% Elementary identities need only 1--3 readable lines, regardless of badge.
% Retain essential hard/required reasoning even when its source badge is SUPPORT.
% Select examples by classroom/learning-check/difficult-concept/counterexample value.
% No one-result-one-example rule: preserve essential cases, delete redundant checks.
% Merge repeated setup, not selected examples into incomplete formula fragments.
% Keep assumptions with statements; use a compact wrapping exam trap only if useful.
% See proof-layout.tex for the optional proof structure.
% Add lookup sheets only for the requested retrieval scope; see retrieval-and-scope.md.
% Keep a notation registry rather than globally replacing overloaded symbols.
% Put source-optional extensions at lecture end or in an authorized separate volume.
% Use canonical proof labels and actual references; delete later-covered previews.
% Separate lookup/self-test volumes may be authorized; avoid duplication in the body.
% Keep standalone body definitions and move experiment minutiae to the source map.
% For compression, baseline main-body pages/words and set a task-specific budget.
% Report main vs total vs relocated; no fixed course quota and no tiny-font workaround.
% Inspect actual page counts, blank space, and orphan Needspace failures.
% Assessed APPLY material stays in the main study path even when its proof is optional.
% Identify computed, constructed, and source-measured values separately.
\renewcommand{\reviewtitle}{Course review}
% A neutral header stays accurate when two short lectures share a page.
\renewcommand{\reviewtopic}{Illustrative notes}
\begin{document}

\reviewsection{Lecture 1: Center and spread}{lecture-one}
\begin{reviewsynthesis}{Concept comparison}
{\arrayrulecolor{ReviewBlue!40}
\begin{tabularx}{\linewidth}{@{}P{29mm}|Y|Y@{}}
\textbf{Dimension} & \textbf{Population variance} & \textbf{Sample variance} \\\hline
Symbol & $\sigma^2$ & $s^2$ \\\hline
Reference center & Population mean $\mu$ & Sample mean $\bar x$ \\\hline
Denominator & $N$ & $n-1$ \\
\end{tabularx}}
\end{reviewsynthesis}

\begin{workedexample}{Calculate both variances}
\given{Values $1,2,3$. Compare population and sample variance, with $N=n=3$.}
\begin{enumerate}
\item \textbf{Mean:}
\[
\bar x=\frac1n\sum_{i=1}^n x_i=\frac{1+2+3}{3}=2.
\]
\item \textbf{Population variance:} for the population interpretation, $\mu=\bar x$.
\[
\sigma^2=\frac1N\sum_{i=1}^N(x_i-\mu)^2
=\frac{(1-2)^2+(2-2)^2+(3-2)^2}{3}
=\frac23\approx0.6667.
\]
\item \textbf{Sample variance:}
\[
s^2=\frac1{n-1}\sum_{i=1}^n(x_i-\bar x)^2
=\frac{(1-2)^2+(2-2)^2+(3-2)^2}{3-1}=1.
\]
\end{enumerate}
\result{The denominator changes from $N$ to $n-1$. Choose it from the question's population/sample interpretation.}
\end{workedexample}

\begin{reviewsynthesis}{Lecture 1 synthesis}
{\arrayrulecolor{ReviewBlue!40}
\begin{tabularx}{\linewidth}{@{}P{29mm}|Y|Y@{}}
\textbf{Operation} & \textbf{Quantity} & \textbf{Relationship} \\\hline
Center & Mean & Reference for deviations \\\hline
Subtract the center & $x_i-\mu$ or $x_i-\bar x$ & Signed deviation in the original units \\\hline
Aggregate squares & Variance & Denominator depends on population or sample \\\hline
Take the square root & Standard deviation & Returns to the original units \\
\end{tabularx}}
\end{reviewsynthesis}

% Do not force a new page for every short lecture; inspect the rendered break.
\reviewsection{Lecture 2: Neighbor voting}{lecture-two}
\begin{reviewsynthesis}{Concept and adjacent instance}
{\arrayrulecolor{ReviewBlue!40}
\begin{tabularx}{\linewidth}{@{}P{29mm}|Y@{}}
\textbf{Parent concept} & Distance-weighted voting \\\hline
Property & Closer neighbors receive more weight. \\\hline
Instance & Inverse-square weighting: $w_i=1/d_i^2$ \\\hline
Condition & Distances must be nonzero for this expression. \\
\end{tabularx}}
\end{reviewsynthesis}

\begin{workedexample}{Compare the voting rules}
% Keep the complete Given and first action together; the style reserves six lines.
% Inspect the rendered opening if longer givens or a taller first action are used.
\given{Three neighbors at distances $1,2,3$, with labels $A,B,B$. Compare equal votes and inverse-square weights.}
\noindent\begin{minipage}[c]{.57\linewidth}\raggedright
\textbf{1. Equal votes:} $B$ has two votes and $A$ has one. Use the same neighbors for the weighted comparison.
\end{minipage}\hfill
\begin{minipage}[c]{.40\linewidth}\centering
\begin{tikzpicture}[x=.95cm,y=.9cm,font=\sffamily\fontsize{8.5}{10}\selectfont]
\draw[ReviewMuted!45] (-2.4,0)--(3.4,0);
\fill[ReviewBlue] (-2,0) circle (2pt) (3,0) circle (2pt);
\fill[ReviewInk] (0,.10)--(.10,0)--(0,-.10)--(-.10,0)--cycle;
\fill[orange!75!black] (1,.12)--(.88,-.08)--(1.12,-.08)--cycle;
\node[above=4pt,text=ReviewBlue] at (-2,0) {$B$};
\node[above=4pt,text=orange!75!black] at (1,0) {$A$};
\node[above=4pt,text=ReviewBlue] at (3,0) {$B$};
\node[below=4pt] at (-2,0) {$d=2$};
\node[above=4pt] at (0,0) {query};
\node[below=4pt] at (1,0) {$d=1$};
\node[below=4pt] at (3,0) {$d=3$};
\node[text=ReviewMuted] at (.5,-1.05) {One possible 1-D layout};
\end{tikzpicture}
\end{minipage}\par
\textbf{2. Inverse-square weights:}
\[
\begin{aligned}
W_A&=\sum_{i:y_i=A}\frac1{d_i^2}
=\frac1{1^2}=1,\\[3pt]
W_B&=\sum_{i:y_i=B}\frac1{d_i^2}
=\frac1{2^2}+\frac1{3^2}=\frac{13}{36}.
\end{aligned}
\]
\result{Equal voting predicts $B$, whereas inverse-square voting predicts $A$.}
\end{workedexample}

\begin{reviewsynthesis}{Lecture 2 synthesis}
{\arrayrulecolor{ReviewBlue!40}
\begin{tabularx}{\linewidth}{@{}P{29mm}|Y|P{22mm}@{}}
\textbf{Rule} & \textbf{Role of distance} & \textbf{This example} \\\hline
Equal votes & Distances select the neighbors. & Class $B$ \\\hline
Inverse-square weights & Distances also determine voting weight. & Class $A$ \\
\end{tabularx}}
\end{reviewsynthesis}
\end{document}
